\documentclass[a4paper,11pt]{article}

\usepackage[utf8]{inputenc}
\usepackage[T1]{fontenc}
\usepackage[ngerman]{babel}
\usepackage{amsmath}
\usepackage{amssymb}
\usepackage{xcolor}
\usepackage{tikz}
\usepackage{enumitem}
\usepackage{needspace}
\usepackage{geometry}
\geometry{a4paper,top=1.4cm,bottom=1.4cm,left=1.8cm,right=1.8cm}

% ----- Dark Mode (kitty-Theme: kitty.conf) -----
\definecolor{bgdark}{HTML}{11121A}   % kitty background
\definecolor{fgtext}{HTML}{D5DDE8}   % kitty foreground
\definecolor{accent}{HTML}{91A8D0}   % kitty color4 (blau) -- Kapitel
\definecolor{accent2}{HTML}{E8D49A}  % kitty color11 (gelb) -- Aufgaben
\definecolor{accent3}{HTML}{B5D4A8}  % kitty color10 (grün) -- Selbstreflexion
\definecolor{accent4}{HTML}{D4BCE6}  % kitty color13 (violett) -- Denkanstoß
\definecolor{rulecol}{HTML}{4A5B78}  % kitty color8
\definecolor{gridcol}{HTML}{3A4A63}  % Karo
\pagecolor{bgdark}
\color{fgtext}

\setlength{\parindent}{0pt}
\setlength{\parskip}{4pt}

% Karo-Raster; #1 = Höhe in cm
\newcommand{\karo}[1]{\par\noindent
  \begin{tikzpicture}
    \draw[step=5mm, gridcol, very thin] (0,0) grid (\linewidth,#1);
  \end{tikzpicture}\par}

\newcommand{\kapitel}[1]{%
  \clearpage
  {\Large\bfseries\color{accent} #1}\par
  \noindent\textcolor{accent}{\rule{\linewidth}{0.8pt}}\par\vspace{4pt}}

% Aufgabenkopf; #1 = ID/Titel, #2 = Typ
\newcommand{\aufg}[2]{%
  \par\vspace{4pt}
  {\bfseries\color{accent2} #1}~~{\footnotesize\color{rulecol}[#2]}\par\vspace{2pt}}

\newcommand{\teaser}[1]{%
  \par\vspace{2pt}{\color{accent4}\textbf{Denkanstoß.}}~#1\par}
\newcommand{\marg}[1]{\par\vspace{1pt}{\footnotesize\itshape\color{rulecol}Randnotiz: #1}\par}
\newcommand{\reflexion}{%
  \par\vspace{2pt}{\bfseries\color{accent3} Selbstreflexionsfragen}\par\vspace{1pt}}

\setlist[itemize]{leftmargin=1.4em,itemsep=1pt,topsep=1pt,parsep=0pt,label=\textcolor{accent3}{\textbullet}}

\begin{document}

% ---------- Titelkopf ----------
{\Large\bfseries\color{accent} Numerische Methoden -- Aufgabenblatt}\\[4pt]
{\small\color{fgtext} Alle Aufgaben aus dem Vorlesungsskript \emph{``Numerical Methods''}
(Prof.\ Dr.-Ing.\ M.\ Schutera) -- 1:1 extrahiert, mit Platz zum Ausfüllen.}\\[10pt]
\noindent\textcolor{rulecol}{\rule{\linewidth}{0.6pt}}\\[6pt]
{\small\color{fgtext}\textbf{Name:} \rule{6cm}{0.3pt} \hfill \textbf{Datum:} \rule{4cm}{0.3pt}}\\[8pt]
\noindent\textcolor{rulecol}{\rule{\linewidth}{0.6pt}}

% #####################################################################
\kapitel{Kapitel 1 \;--\; Einstieg in die Numerischen Methoden}

\needspace{19cm}
\aufg{1.1 \;--\; Minimum der Sinusfunktion}{Handrechnung}
Bestimmen Sie das Minimum von $\sin(x)$ auf $[-3,1]$, indem Sie die Funktion mit
Schrittweite $h$ diskretisieren, an den Gitterpunkten auswerten und den kleinsten
Wert auswählen. Vergleichen Sie mit dem analytisch exakten Minimum.
\karo{16}

\needspace{20cm}
\aufg{1.2 \;--\; Grid-Search}{Handrechnung}
Führen Sie eine Grid-Search für $w,b$ mit Schrittweite $2{,}5$ über
$\{0;2{,}5;5;7{,}5;10\}$ durch (System
$\big(\begin{smallmatrix}2&1\\5&1\end{smallmatrix}\big)\big(\begin{smallmatrix}w\\b\end{smallmatrix}\big)=\big(\begin{smallmatrix}11\\13\end{smallmatrix}\big)$).
Fehler $\mathrm{error}_1=|2w+b-11|$, $\mathrm{error}_2=|5w+b-13|$. Wählen Sie das
Paar $(w,b)$ mit kleinstem akkumuliertem Fehler.
\karo{16}

\needspace{11cm}
\aufg{1.3 \;--\; Enter the machine}{Code}
Setzen Sie die Grid-Search im vorgehosteten Jupyter-Notebook (Python-Vorlage) um.
Nutzen Sie die Selbstreflexionsfragen als Leitfaden.
\karo{7}

\needspace{8cm}
\teaser{Fällt Ihnen eine einfache Möglichkeit ein, das Verhältnis von Genauigkeit zu
Rechenaufwand der Grid-Search zu verbessern?}
\karo{5}

\needspace{12cm}
\reflexion
\begin{itemize}
  \item Warum ist die Wahl der Schrittweite $h$ wichtig (Genauigkeit vs.\ Rechenzeit)?
  \item Wie beeinflusst das Intervall $[a,b]$ Ergebnisse und Lage der Extrema?
  \item Unterschiede stetige vs.\ diskretisierte Funktion -- Implikationen?
\end{itemize}
\karo{8}

% #####################################################################
\kapitel{Kapitel 2 \;--\; Gleitkomma-Arithmetik}

\needspace{20cm}
\aufg{2.1 \;--\; Rundungsfehler}{Handrechnung}
Betrachten Sie $10/3$. Schreiben Sie die Dezimalentwicklung aus und bestimmen Sie den
Rundungsfehler beim Abschneiden nach drei Nachkommastellen.
\karo{17}

\needspace{19cm}
\aufg{2.2 \;--\; 2-Bit-Mantisse \& Maschinenepsilon}{Handrechnung}
Eine unnormalisierte Mantisse hat $2$ Bit ($0{.}xx_2$). Bestimmen Sie alle
darstellbaren Werte und $\varepsilon_{\mathrm{mach}}=2^{-t}$. Relative vs.\ absolute
Präzision?
\karo{16}

\needspace{20cm}
\aufg{2.3 \;--\; Auslöschung von Hand}{Handrechnung}
Zwei Zahlen, $8$ signifikante Stellen.
\textbf{(a)} $a=12345678{,}5$, $b=12345678{,}0$: $\tilde a,\tilde b,\ \tilde a-\tilde b$
vs.\ $a-b$.
\textbf{(b)} $a=12345678{,}4$, $b=12345678{,}0$: ebenso. Beobachtung beim relativen
Fehler?
\karo{15}

\needspace{18cm}
\aufg{2.4 \;--\; Festkomma-Skalierung}{Handrechnung}
Bereich $[-1000,1000]$, 32-Bit-Integer, Skalierungsfaktor $10^{6}$. Wie wird
$1{,}234567$ gespeichert? Präzision (kleinste Differenz) und maximaler Rundungsfehler?
\karo{16}

\needspace{12cm}
\aufg{2.5 \;--\; Maschinenepsilon (hands on)}{Code}
Bestimmen Sie das Maschinenepsilon eines Systems per Pseudocode. Festkomma vs.\
Gleitkomma? Welcher Datentyp für einen Taschenrechner und warum?
\karo{8}

\needspace{10cm}
\aufg{2.6 \;--\; Auslöschung demonstrieren (hands on)}{Code}
Wie demonstrieren Sie Auslöschung am Rechner? Schreiben Sie Pseudocode auf, bevor Sie
weiterlesen.
\karo{7}

\needspace{13cm}
\aufg{2.7 \;--\; Argumentation}{Konzept}
Wie leistet $f(\epsilon)=\dfrac{1}{a+\epsilon-b}$ für kleines $\epsilon$ und $a=b$
dies? Was erwarten Sie für sehr kleines $\epsilon$?
\marg{Was passiert für $a>b$? Denken Sie in signifikanten Stellen.}
\karo{8}

\needspace{8cm}
\teaser{Fallen Ihnen Metriken für numerische Verfahren ein, basierend auf den
besprochenen Approximationen und Fehlern?}
\karo{5}

\needspace{12cm}
\reflexion
\begin{itemize}
  \item Aufbau und Komponenten der Gleitkommadarstellung?
  \item Was ist das Maschinenepsilon, Bezug zur Genauigkeit?
  \item Praktische Auswirkungen auf numerische Berechnungen?
\end{itemize}
\karo{8}

% #####################################################################
\kapitel{Kapitel 3 \;--\; Fehleranalyse}

\needspace{22cm}
\aufg{3.1 \;--\; Minimum der Sinusfunktion (Eigenschaften)}{Handrechnung}
Minimum von $\sin(x)$ auf $[-3,1]$ erneut; diskretisiert mit $h=1$ und $h=0{,}2$;
Eingabe gestört $\tilde x=x\pm0{,}25$. Analysieren und quantifizieren Sie
\textbf{Kondition}, \textbf{Stabilität}, \textbf{Konsistenz} und \textbf{Konvergenz}
in beiden Fällen mit den Formeln aus dem vorigen Abschnitt.
\karo{18}

\needspace{20cm}
\aufg{3.2 \;--\; Stabilität und Kondition}{Beweis}
Zeigen Sie: Die Stabilität entspricht der Kondition des Systems für $h\to0$.
\karo{17}

\needspace{12cm}
\aufg{3.3 \;--\; Compute-getriebene Fehleranalyse (hands on)}{Code}
Brute-forcen Sie die Fehleranalyse beim Lösen des Zwei-Gleichungs-Systems aus Kap.\,1;
bestimmen und optimieren Sie Kondition/Stabilität/Konsistenz/Konvergenz.
\karo{8}

\needspace{14cm}
\reflexion
\begin{itemize}
  \item Konditionierung des Zwei-Gleichungs-Systems (vgl.\ $\min\sin$)?
  \item Konvergiert die Stabilität mit kleineren Schrittweiten? Wogegen?
  \item Wirkung von Eingabestörungen; Zusammenspiel mit $h$?
  \item Konvergenz vs.\ Stabilität \& Konsistenz (Bias/Varianz)?
  \item Genauigkeitsgrenze für immer kleinere $h$? Warum?
  \item Welches weitere Problem beim Verfeinern von $h$?
\end{itemize}
\karo{8}

\needspace{8cm}
\teaser{Fällt Ihnen eine Möglichkeit ein, Brute-Force zu verfeinern, um mit weniger
Aufwand bessere Ergebnisse zu erzielen?}
\karo{5}

% #####################################################################
\kapitel{Kapitel 4 \;--\; Newton-Verfahren}

\needspace{20cm}
\aufg{4.1 \;--\; Newton für $\sqrt{2}$}{Handrechnung}
Berechnen Sie $\sqrt{2}$ mit Newton ($f(x)=x^2-2$), Start $x_0=2$. Einige
Iterationen; beobachten Sie das Quadrieren des Fehlers.
\karo{16}

\needspace{20cm}
\aufg{4.2 \;--\; Newton vs.\ Grid-Search für $\sin(x)=0$}{Handrechnung}
Nullstelle von $\sin(x)=0$ auf $[2{,}5;4]$ (nahe $\pi$).
\textbf{(a)} Grid-Search, $h=0{,}3$. \textbf{(b)} Newton ($f'=\cos$), $x_0=3{,}5$.
\textbf{(c)} Genauigkeit \& Auswertungen vergleichen.
\karo{15}

\needspace{18cm}
\aufg{4.3 \;--\; Denksportaufgabe: Versagensmodi}{Schaubild}
Visualisieren Sie die Versagensmodi des Newton-Verfahrens im Schaubild von
$f(x)=x^3-x$.
\karo{15}

\needspace{20cm}
\aufg{4.4 \;--\; Newton von Hand}{Handrechnung}
Nullstelle von $f(x)=x^3-x=0$ (Nullstellen $-1,0,1$), erst geometrisch, dann von Hand;
gesucht $x=1$ ab $x_0=1{,}4$.
\karo{16}

\needspace{20cm}
\aufg{4.5 \;--\; Sekantenverfahren von Hand}{Handrechnung}
Sekantenverfahren auf $f(x)=x^3-x$ mit $x_0=1{,}2$, $x_1=1{,}4$; erst geometrisch,
dann von Hand.
\karo{16}

\needspace{19cm}
\aufg{4.6 \;--\; Geometrische Fehlersuche}{Handrechnung}
Startpunkte $x_0$ für drei Versagensmodi: verschwindende Ableitung; Konvergenz gegen
eine Nicht-Ziel-Nullstelle; Oszillation zwischen Punkten.
\karo{16}

\needspace{11cm}
\aufg{4.7 \;--\; Enter the machine}{Code}
Konvergenz(-raten) von Grid-Search, Newton, Sekanten in Python; Versagensmodi über
$x_0$ erkunden (geometrisch, analytisch, im Code).
\karo{7}

\needspace{12cm}
\reflexion
\begin{itemize}
  \item Hauptvorteil von Newton gegenüber Grid-Search?
  \item Warum genügt lineare Approximation (Taylor)?
  \item Mächtigkeit quadratischer Konvergenz ($10^{-1}\!\to\!10^{-16}$?).
  \item Wann Sekanten vorziehen, wann nicht?
\end{itemize}
\karo{8}

\needspace{8cm}
\teaser{Wie stellen wir sicher, dass wir die globale Lösung finden, nicht nur eine
lokale?}
\karo{5}

% #####################################################################
\kapitel{Kapitel 5 \;--\; Globale Optimierung}

\needspace{11cm}
\aufg{5.1 \;--\; Newton auf Shekel, weiterer Startpunkt}{Code}
Newton auf der 1D-Shekel-Funktion von verschiedenen Startpunkten; wohin konvergiert es?
Diskutieren Sie das allgemeine Verhalten.
\marg{Wie ändern Sie die Update-Regel, um stattdessen Maxima zu finden?}
\karo{7}

\needspace{20cm}
\aufg{5.2 \;--\; Random Restarts}{Handrechnung}
$5$ lokale Minima, globales Becken $p=0{,}15$. $P$ bei $K=10$ ($P=1-(1-p)^K$); $K$ für
$99\%$ ($K=\frac{\ln(1-P)}{\ln(1-p)}$); dann $K$ für $p=0{,}05$ und $p=0{,}01$.
\karo{15}

\needspace{11cm}
\aufg{5.3 \;--\; Eigene 1D-Shekel-Funktion}{Code}
Eigene 1D-Shekel-Funktion entwerfen; lokale Verfahren anwenden, Newton-Versagen zeigen;
verschiedene $x_0$ untersuchen.
\karo{7}

\needspace{11cm}
\aufg{5.4 \;--\; Die Suche zu Minima lenken}{Konzept + Code}
$|f''|$-Trick: Vorzeichen der Krümmung im Nenner umkehren $\to$ Schritt geht stets
bergab. Geometrische Bedeutung? Dann im Code testen.
\karo{8}

\needspace{12cm}
\aufg{5.5 \;--\; Einzugsgebiet (Basin of Attraction)}{Schaubild}
Markieren Sie in den Shekel-Plots die Einzugsgebiete jedes lokalen Minimums und
schätzen Sie ihre relativen Größen. Welches ist global? Bezug zur Random-Restart-Wahrscheinlichkeit?
\karo{9}

\needspace{11cm}
\aufg{5.6 \;--\; Lokalen Minima entkommen}{Code}
Random Restarts und Basin Hopping auf Ihrer Shekel-Funktion; Funktionsiterationen bis
zum globalen Minimum vergleichen. Einfluss der Schrittweitenverteilung?
\karo{7}

\needspace{9cm}
\aufg{5.7 \;--\; Nicht-konvexe Becken}{Konzept}
Notieren Sie eine Funktion, bei der die Einzugsgebiete nicht leicht abzugrenzen sind.
\karo{6}

\needspace{10cm}
\aufg{5.8 \;--\; Adaptive Schrittweite}{Konzept + Code}
Geschickter Weg, die Schrittweitenverteilung automatisch anzupassen? Gute Heuristik?
Wie katastrophale Sprünge verhindern? Implementieren.
\karo{7}

\needspace{9cm}
\aufg{5.9 \;--\; Noise and landscape engineering}{Konzept}
Wie wirkt das Approximieren der Loss-Landschaft auf Mini-Batches (unterschiedliche
Punktauswahl)?
\karo{6}

\needspace{10cm}
\aufg{5.10 \;--\; Code-Übung: Restarts vs.\ Basin Hopping}{Code}
Random Restarts und Basin Hopping auf der 1D-Shekel-Funktion implementieren;
Funktionsauswertungen bis zum globalen Minimum bei $x\approx4$ vergleichen.
\karo{7}

\needspace{20cm}
\aufg{5.11 \;--\; Einzugsgebiete kartieren}{Handrechnung}
Lokaler Optimierer $\to$ nächstes Foxhole.
\textbf{(a)} Einzugsgebiete der Foxholes bei $x=0,4,7$ skizzieren -- Basin-Grenzen?
\textbf{(b)} Sprung $\delta\sim\mathrm{Uniform}(-3,3)$ ab $x^{\circ}=7$: $P$(im
globalen Becken)?
\textbf{(c)} Warum Basin Hopping hier evtl.\ schneller als Random Restarts?
\karo{14}

\needspace{14cm}
\reflexion
\begin{itemize}
  \item Warum versagt lokale Optimierung auf nicht-konvexen Landschaften?
  \item Aufwand vs.\ Größe des globalen Beckens (Random Restarts)?
  \item Basin Hopping vs.\ Random Restarts?
  \item Was sagt $f''(x_n)$ über die Loss-Landschaft, wie nutzbar machen?
  \item Wie helfen verrauschte (Mini-Batch-)Schätzungen beim Entkommen?
  \item $x_0$ in einem lokalen Becken: Wege zum Entkommen -- Wirksamkeit?
\end{itemize}
\karo{8}

\needspace{8cm}
\teaser{Was, wenn wir einer hochfrequenten Version von Shekels Foxholes
gegenüberstehen? Welches Problem tritt auf?}
\karo{5}

% #####################################################################
\kapitel{Kapitel 6 \;--\; Numerische Integration}

\needspace{18cm}
\aufg{6.1 \;--\; Exakter Wert}{Handrechnung}
Berechnen Sie $\int_{-2}^{-1}\sin(x)\,dx$ exakt (Referenzwert).
\karo{16}

\needspace{18cm}
\aufg{6.2 \;--\; Mittelpunktsregel von Hand}{Handrechnung}
$\int_{-2}^{-1}\sin(x)\,dx$ mit der Mittelpunktsregel, $n=1$.
\karo{16}

\needspace{18cm}
\aufg{6.3 \;--\; Mittelpunktsregel mit mehr Intervallen}{Handrechnung}
Mittelpunktsregel mit $n=2$ (Mittelpunkte $-1{,}75$ und $-1{,}25$).
\marg{Berechnen Sie die Approximation auch an den Endpunkten und vergleichen Sie.}
\karo{15}

\needspace{18cm}
\aufg{6.4 \;--\; Trapezregel von Hand}{Handrechnung}
Trapezregel, $n=2$; Stützstellen $-2,-1{,}5,-1$.
\karo{16}

\needspace{18cm}
\aufg{6.5 \;--\; Simpson-Regel von Hand}{Handrechnung}
Simpson-Regel, $n=2$; Stützstellen $-2,-1{,}5,-1$.
\marg{Mehr Intervalle; zusammengesetzte Simpson-/Trapezregeln.}
\karo{15}

\needspace{18cm}
\aufg{6.6 \;--\; Monte Carlo von Hand}{Handrechnung}
Monte Carlo, $N=4$, Stichproben $X_1=-1{,}95$, $X_2=-1{,}55$, $X_3=-1{,}15$,
$X_4=-1{,}05$.
\karo{15}

\needspace{19cm}
\aufg{6.7 \;--\; Monte Carlo: Fehler \& $N$ erhöhen}{Konzept}
Stichprobenvarianz und Standardfehler bestimmen. Was passiert bei $N=16$? Warum ist
der MC-Fehler dimensionsunabhängig?
\karo{15}

\needspace{20cm}
\aufg{6.8 \;--\; Simpson-Gewichte via Lagrange}{Herleitung}
Leiten Sie die Simpson-Gewichte her, indem Sie die Lagrange-Basispolynome über
$[a,b]$ integrieren.
\marg{Beginnen Sie mit $x=0$; was passiert, wenn Sie eine lineare Funktion quadratisch
approximieren?}
\karo{15}

\needspace{11cm}
\aufg{6.9 \;--\; Höhere Dimensionen (on your machine)}{Code}
MC-Integration für ein $d$-dimensionales Integral (multivariate Gauß über
Hyperwürfel); Fehler über $d$; Versagen der Gitter-Quadratur.
\karo{7}

\needspace{16cm}
\reflexion
\begin{itemize}
  \item Fehlervergleich Mittelpunkt/Trapez/Simpson/Monte Carlo?
  \item Warum Mittelpunkt besser als Links-/Rechts-Endpunkt?
  \item Zusammengesetzte Methoden: Genauigkeit \& Trade-off?
  \item Warum versagt deterministische Quadratur in hohen Dimensionen?
  \item Mittelpunkt als Spezialfall von Monte Carlo?
  \item Intuition: warum ``explodiert'' MC in hohen Dim.\ nicht?
  \item Wann Simpson, wann Monte Carlo?
  \item Mini-Batch-Mittelwert in SGD und Monte Carlo?
  \item Wie hilft die Gradientenschätzung in SGD beim Entkommen?
\end{itemize}
\karo{8}

\needspace{8cm}
\teaser{Warum werden hochgradige Quadraturmethoden instabil, wenn sie Polynome
niedrigen Grades approximieren?}
\karo{5}

% #####################################################################
\kapitel{Kapitel 7 \;--\; Modelltraining}

\needspace{20cm}
\aufg{7.1 \;--\; Loss bei Initialisierung}{Handrechnung}
$\hat y=wx+b$, $L=\tfrac1n\sum_i(wx_i+b-y_i)^2$, Init $w=0,\ b=\bar y$. Zeigen Sie
$L(0,\bar y)=\mathrm{Var}(y)$. Mit $\bar y=231/14=16{,}5$, $\sum(y_i-\bar y)^2=1401{,}5$:
erster Loss-Wert? Bedeutung von $5000$ bzw.\ $0{,}0001$?
\karo{15}

\needspace{13cm}
\aufg{7.2 \;--\; Overfitting als Sanity-Check}{Konzept}
Grad-4-Polynom durch zwei Punkte ($(4{,}0;4)$, $(9{,}5;21)$). Wie viele Null-Loss-Kurven
gehen durch beide Punkte, und was sagt das über die Geometrie der Verlustlandschaft?
\karo{9}

\needspace{12cm}
\aufg{7.3 \;--\; Mini-Batch-SGD}{Konzept}
In welchem Sinn ist Mini-Batch-SGD ein Monte-Carlo-Schätzer (was wird worüber
gemittelt)? Ab welchem $n$ kippt der Engpass, gegeben Fehler $\propto 1/\sqrt{B}$?
\karo{8}

\needspace{11cm}
\aufg{7.4 \;--\; Mehrere Random Seeds}{Konzept}
$\hat y=w^2x+b$: warum ist der Ursprung ein Sattel, warum ist Null-Init eine Falle?
Wozu Training aus mehreren Random Seeds?
\karo{8}

\needspace{11cm}
\aufg{7.5 \;--\; Quantisierung für Inferenz}{Konzept}
Warum versendet man Modelle in float32, wenn int hier fast gleich gut war? Warum lässt
sich hochpräzise trainiert niedrigpräzise einsetzen, während Training niedrigpräzise
scheitert?
\karo{8}

\needspace{12cm}
\reflexion
\begin{itemize}
  \item Warum entdeckt Loss$\to0$ bei Überparametrisierung Pipeline-Bugs statt
        Optimierer-Bugs?
  \item Was sagt ein erster Loss von $5000$ bzw.\ $0{,}0001$ (statt $\approx100$)?
  \item Mini-Batch-SGD als Monte-Carlo-Schätzer -- in welchem Sinn?
  \item $\hat y=w^2x+b$: Ursprung als Sattel, Null-Init als Falle?
  \item Hochpräzise trainieren, niedrigpräzise einsetzen -- warum geht das?
\end{itemize}
\karo{8}

\end{document}
